\documentclass[11pt]{article}
\usepackage{icmlko}
\begin{document}

\title{제일 큰 시장에서만 진다 --- 덱에 두 해 동안 걸려 있던 배수를 다시 재니 순서가 뒤집혔다}
\author{월드모델 실험실}
\date{노트 580 · 2026년 8월 5일}
\maketitle

\begin{abstract}
노트 $556$ 이 투자자 덱의 숫자 열여섯 곳을 전수 정정하면서 \textbf{하나를
남겼다} --- ``통합 모델 $\div$ 그 시장 전용 모델''(만화 $2.00$ 배 등)은
\emph{노트 $200$ 번대의 다른 챔피언 세대} 수치라 별도 실험이 필요하다고
적어 두었다. 그것을 쟀다.

\textbf{견줌부터 어긋나 있었다.} 옛 수치는 통합 $=$ 유보 $\rho$ 인데 전용
$=$ \emph{교차검증} 최고점이었다 --- 서로 다른 자다. \textbf{둘 다 같은
$2025$ 년 이후 유보에서} 다시 쟀다.

\begin{center}
\begin{tabular}{lrrr}
\hline
 & 덱(옛 세대) & \textbf{이 세대} & \\
\hline
만화 & $2.00$ & $\mathbf{1.03}$ & \\
아이돌 & $1.91$ & $1.36$ & \\
펀딩 & $1.72$ & $1.08$ & \\
모바일 앱 & $1.11$ & $\mathbf{0.98}$ & \textbf{유일하게 진다} \\
도서 & $0.95$ & $\mathbf{3.29}$ & \\
시장 팝업 & --- & $\mathbf{13.85}$ & \\
\hline
\end{tabular}
\end{center}

\textbf{순서가 통째로 바뀐다.} 그런데 \emph{이야기는 더 좋아진다} ---
이득의 크기와 그 시장의 학습 행수가 순위상관 $\mathbf{-0.697}$ 로
맞물린다.
\end{abstract}

\begin{figure}[t]
\centering
\includegraphics[width=\textwidth]{figs/f1.pdf}
\caption{\textbf{가} --- 학습 행이 적을수록 크게 이긴다.
\textbf{나} --- 순서가 뒤집힌다. \textbf{다} --- 덱을 고쳤다.}
\end{figure}

\section{전부}

\begin{center}
\begin{tabular}{lrrrrr}
\hline
시장 & 학습 & 유보 & 통합 & 전용 & 통합$\div$전용 \\
\hline
시장 팝업 & $123$ & $126$ & $+0.1564$ & $+0.0113$ & $\mathbf{13.85}$ \\
도서 & $80$ & $163$ & $+0.4038$ & $+0.1227$ & $\mathbf{3.29}$ \\
아이돌 & $122$ & $51$ & $+0.4686$ & $+0.3458$ & $1.36$ \\
게임 & $259$ & $180$ & $+0.6644$ & $+0.5884$ & $1.13$ \\
애니 & $2{,}779$ & $606$ & $+0.6024$ & $+0.5454$ & $1.10$ \\
세계애니 & $2{,}648$ & $300$ & $+0.5769$ & $+0.5231$ & $1.10$ \\
웹툰 & $2{,}938$ & $650$ & $+0.3563$ & $+0.3283$ & $1.09$ \\
펀딩 & $2{,}358$ & $529$ & $+0.3678$ & $+0.3407$ & $1.08$ \\
만화 & $5{,}647$ & $258$ & $+0.3375$ & $+0.3273$ & $1.03$ \\
\textbf{모바일 앱} & $1{,}559$ & $441$ & $+0.6262$ & $+0.6390$ & $\mathbf{0.98}$ \\
\textbf{팝업} & $\mathbf{24}$ & $65$ & $+0.3049$ & \textbf{못 만든다} & --- \\
\hline
\end{tabular}
\end{center}

\textbf{유보 가중 평균 $\mathbf{1.67}$ 배.} 사전 등록 셋이 다 섰다 ---
① 통합이 이기는 시장 $9/11$ ② 학습 행이 적을수록 이득이 크다($-0.697$)
③ 모바일이 제일 작다($0.98$, 유일하게 진다).

\section{팝업은 전용 모형을 만들 수가 없다}

학습이 $24$ 행이라 $5$겹이 안 선다. \textbf{그것 자체가 이 판이 파는
것이다} --- 새 시장에는 전용 모형을 만들 자료가 없고, 통합은 그래도
$+0.3049$ 를 낸다. 표에서 ``통합이 이기는 $9/11$''은 실은
\textbf{전용을 만들 수 있는 $10$ 중 $9$} 다.

노트 $578$ 이 하루 전에 같은 도메인을 반대편에서 봤다 --- 팝업 $16$ 행을
학습에서 빼면 팝업 자신은 $+0.0461$ 오르는데 판 전체는 안 오른다.
\textbf{팝업은 순 기부자이면서 동시에 최대 수혜자다.} 자기 $16$ 행으로는
못 배우고, 남의 $18{,}000$ 행으로 $+0.30$ 을 얻는다.

\section{옛 수치는 왜 그렇게 달랐나}

두 가지가 겹쳤다.

\begin{itemize}
\item \textbf{자가 달랐다} --- 전용을 교차검증 최고점으로 잰 것은
\emph{그 시장에 유리한} 자다. 그런데도 만화 $2.00$ 배가 나왔다면 그때
전용이 지금보다 훨씬 나빴다는 뜻이다.
\item \textbf{세대가 달랐다} --- 그 사이 판이 $0.4908 \to 0.4663$ 으로
두 번 정정됐고(노트 $547 \cdot 552$), 축이 늘고 행 무게가 바뀌었다.
\end{itemize}

\textbf{어느 쪽이 얼마인지는 안 갈랐다.} 갈라도 덱에 쓸 숫자는 새 것이므로
값이 없다 --- \textbf{안 잰 것을 안 쟀다고 적는다.}

\section{덱에서 지워야 했던 문장 하나}

슬라이드 $11$ 이 ``\textbf{만화는 다른 시장의 도움을 가장 많이 받아
갑니다($2.00$ 배)}''를 적고 있었다. 만화는 $1.03$ 배로 열 시장 중 거의
꼴찌다. 고쳤다 ---

\begin{quote}
만화는 학습의 $30\%$($5{,}647$ 행)를 차지하는데 절반을 덜면 성능이
\emph{오르고}(\texttt{TRAINCAP}$=0.5$, 노트 $509$), 통합에서 받아 가는
것도 $1.03$ 배로 거의 없다. \textbf{제일 많이 넣는데 주지도 받지도 않는
시장}이다.
\end{quote}

\section*{조항}

\begin{quote}
\textbf{두 모형을 견줄 때는 채점하는 표본과 채점하는 자를 먼저 맞춘다.}
한쪽은 유보 $\rho$, 다른 쪽은 교차검증 최고점이면 그 비는 아무것도 재지
않는다. \textbf{같은 조항을 오늘 두 번 적는다} --- 노트 $579$ 는
\emph{축}이 어긋나 있었고 여기서는 \emph{표본}이 어긋나 있었다.
\end{quote}

\section*{남은 것}

\begin{itemize}
\item \textbf{시장 팝업 $13.85$} --- 전용이 $+0.0113$ 이라 비가 불안정하다.
차($+0.1451$)로 적는 편이 정직하다.
\item \textbf{전용 손잡이 격자가 넷뿐이다} --- 더 넓히면 전용이 오른다.
지금 수치는 통합에 유리한 쪽으로 치우쳐 있을 수 있다.
\item \textbf{통합에만 있는 것 둘} --- \texttt{SEASON}$\cdot$\texttt{SPEC}
칸과 행 무게(\texttt{TRAINW}$\cdot$\texttt{TRAINR}). 전용은 그 시장이
관측하는 축칸을 \emph{전부} 받으므로, 이득의 출처는 축이 아니라
\textbf{다른 시장의 행}이다.
\end{itemize}

\end{document}
